\documentclass[11pt,a4paper]{article}

% ── Packages ──────────────────────────────────────────────────────────────────
\usepackage[margin=2.5cm]{geometry}
\usepackage{amsmath,amssymb}
\usepackage{booktabs}
\usepackage{hyperref}
\usepackage{xcolor}
\usepackage{longtable}
\usepackage{array}
\usepackage{graphicx}
\usepackage{fancyhdr}
\usepackage{titlesec}
\usepackage{enumitem}

% ── Colour definitions ────────────────────────────────────────────────────────
\definecolor{apexblue}{RGB}{0,51,102}
\definecolor{apexgray}{RGB}{100,100,100}
\definecolor{apexred}{RGB}{180,30,30}

% ── Section formatting ────────────────────────────────────────────────────────
\titleformat{\section}{\large\bfseries\color{apexblue}}{}{0em}{\thesection\quad}
\titleformat{\subsection}{\normalsize\bfseries\color{apexblue}}{}{0em}{\thesubsection\quad}

% ── Header / footer ───────────────────────────────────────────────────────────
\pagestyle{fancy}
\fancyhf{}
\rhead{\textcolor{apexgray}{\small RESTRICTED --- Model Risk Sensitive}}
\lhead{\textcolor{apexgray}{\small APEX-MDL-0011 \textbar{} FVA MDD v1.0}}
\rfoot{\textcolor{apexgray}{\small \thepage}}
\lfoot{\textcolor{apexgray}{\small Apex Global Bank \textbar{} Quantitative Research Division}}
\renewcommand{\headrulewidth}{0.4pt}
\renewcommand{\footrulewidth}{0.4pt}

% ── Hyperref setup ────────────────────────────────────────────────────────────
\hypersetup{
  colorlinks=true,
  linkcolor=apexblue,
  urlcolor=apexblue,
  pdftitle={Apex Global Bank -- FVA Model Development Document},
  pdfauthor={XVA Quant Team}
}

% ══════════════════════════════════════════════════════════════════════════════
\begin{document}
% ══════════════════════════════════════════════════════════════════════════════

\begin{titlepage}
  \begin{center}
    \vspace*{2cm}
    {\color{apexblue}\rule{\linewidth}{2pt}}\\[0.5cm]
    {\Huge\bfseries\color{apexblue} Model Development Document}\\[0.4cm]
    {\LARGE\bfseries Funding Valuation Adjustment (FVA)}\\[0.3cm]
    {\Large Version 1.0}\\[0.5cm]
    {\color{apexblue}\rule{\linewidth}{2pt}}\\[1cm]
    {\large Apex Global Bank \textbar{} Quantitative Research Division}\\[0.3cm]
    {\large\textcolor{apexred}{Classification: RESTRICTED --- Model Risk Sensitive}}\\[0.3cm]
    {\large SR 11-7 Section Reference: Model Development}\\[2cm]
    \begin{tabular}{ll}
      \textbf{Model ID:}   & APEX-MDL-0011 \\[4pt]
      \textbf{Date:}       & March 2026 \\[4pt]
      \textbf{Status:}     & In Validation \\[4pt]
      \textbf{Risk Tier:}  & Tier 1 \\
    \end{tabular}
  \end{center}
\end{titlepage}

% ── Table of Contents ─────────────────────────────────────────────────────────
\tableofcontents
\newpage

% ══════════════════════════════════════════════════════════════════════════════
\section{Document Metadata}
% ══════════════════════════════════════════════════════════════════════════════

\begin{table}[h!]
\centering
\caption{Document Control Information}
\begin{tabular}{@{}ll@{}}
\toprule
\textbf{Field} & \textbf{Value} \\
\midrule
Model Name            & Funding Valuation Adjustment (FVA) \\
Model ID              & APEX-MDL-0011 \\
Version               & 1.0 \\
Status                & In Validation \\
Risk Tier             & Tier 1 \\
Regulatory Framework  & IFRS 13 Fair Value; Internal FTP Policy; Basel III (informational) \\
Date of Development   & March 2026 \\
Business Owner        & Head of Global Markets Trading (James Okafor) \\
Model Developer       & Quantitative Research --- XVA Team \\
MRO Review            & Pending --- Dr.\ Rebecca Chen \\
CRO Approval          & Pending --- Dr.\ Priya Nair \\
\bottomrule
\end{tabular}
\end{table}

\subsection{Version History}

\begin{table}[h!]
\centering
\caption{Document Version History}
\begin{tabular}{@{}lll@{}}
\toprule
\textbf{Version} & \textbf{Date} & \textbf{Change Summary} \\
\midrule
0.1 & Jan 2026 & Initial framework --- funding cost only (FCA) \\
0.8 & Feb 2026 & Added funding benefit (FBA); asymmetric treatment \\
1.0 & Mar 2026 & Production candidate; submitted for validation \\
\bottomrule
\end{tabular}
\end{table}

% ══════════════════════════════════════════════════════════════════════════════
\section{Executive Summary}
% ══════════════════════════════════════════════════════════════════════════════

\subsection{Model Purpose}

FVA captures the cost (or benefit) of funding derivatives positions that are not fully
collateralised under a CSA. When Apex enters an uncollateralised derivative with a client
and hedges it in the interbank market under a collateralised CSA, the bank faces an asymmetric
funding situation: variation margin must be posted to the hedge counterparty when the hedge is
out-of-the-money, but no corresponding margin is received from the client. This funding gap
must be financed at Apex's unsecured borrowing rate. FVA quantifies the present value of this
cost over the life of the trade.

\[
\mathrm{FVA} = \mathrm{FCA} - \mathrm{FBA}
\tag{1}
\]

FCA (Funding Cost Adjustment) is the cost of funding receivables when Apex is owed money but
cannot receive collateral. FBA (Funding Benefit Adjustment) is the benefit when Apex owes money
and can use its own funds rather than paying market rates.

\subsection{The FVA Debate}

FVA is the most debated valuation adjustment in derivatives pricing. Hull and White (2012)
argued that FVA adjustments represent double-counting with DVA: a bank's funding costs are
already reflected in its credit spread, and applying FVA separately misprices derivatives by
incorporating an avoidable cost into an asset's fair value. The counterargument --- adopted by
most major dealers --- is that FVA is real: it affects actual cash flows and trading decisions,
regardless of theoretical arguments about perfect markets.

Apex's policy (approved by the CFO and CRO in 2024) is to \textbf{include FVA in derivatives
fair value} for internal pricing and P\&L purposes, consistent with industry practice. This
document acknowledges the theoretical debate and documents the specific methodological choices
made.

\subsection{Output Description}

\begin{itemize}[leftmargin=*]
  \item \textbf{FVA charge per new trade}: Used by front office to price uncollateralised client derivatives
  \item \textbf{Portfolio FVA}: Daily P\&L impact reported in income statement
  \item \textbf{FVA Greeks}: Sensitivities to interest rates (DV01), FX, and credit spreads --- used for hedging
  \item \textbf{FCA/FBA decomposition}: Separated for internal cost allocation between Treasury (funding) and trading (derivative)
  \item \textbf{FVA by counterparty and desk}: Granular attribution for profitability analysis
\end{itemize}

\subsection{Key Assumptions}

\begin{enumerate}[leftmargin=*]
  \item Apex's unsecured borrowing spread over OIS/SOFR is a reliable proxy for the marginal cost of funding derivatives positions.
  \item The funding curve is flat across tenors within a 1-year horizon (curve shape used for longer tenors).
  \item The hedge for every uncollateralised client trade is assumed to be executed collateralised (full VM under ISDA SIMM).
  \item DVA and FVA are computed independently and do not overlap.
  \item Funding costs are symmetric: the same spread applies to borrowing and investing (relaxed in the FCA/FBA asymmetric variant).
\end{enumerate}

\subsection{Known Limitations}

\begin{itemize}[leftmargin=*]
  \item The assumption of a single funding curve ignores liquidity segmentation (different tenors trade at different spreads, especially in stress).
  \item FVA is sensitive to the assumed hedging strategy; actual hedging behaviour may differ from the ``always hedge collateralised'' assumption.
  \item FBA recognition is conservative --- many practitioners argue FBA should not be recognised as it relies on Apex defaulting.
  \item The model does not capture dynamic funding strategies (e.g., using repo to fund derivatives positions more cheaply).
\end{itemize}

% ══════════════════════════════════════════════════════════════════════════════
\section{Business Context and Purpose}
% ══════════════════════════════════════════════════════════════════════════════

\subsection{The Funding Gap Problem}

Consider a simple example:
\begin{enumerate}[leftmargin=*]
  \item Apex writes a 5-year interest rate swap to a corporate client (uncollateralised --- no CSA in place).
  \item Apex hedges by entering the opposite swap with a dealer (collateralised --- daily VM under ISDA SIMM).
  \item If rates move so that the dealer hedge is out-of-the-money for Apex, Apex must post variation margin to the dealer immediately.
  \item But Apex cannot call margin from the corporate client.
  \item Apex must borrow funds to post that margin --- at its unsecured borrowing spread over OIS.
\end{enumerate}

This is the funding gap. Over a 5-year swap, the expected cumulative cost of this funding gap
--- the FVA --- can represent a significant fraction of the swap's gross profit.

\subsection{Market Practice}

Following the 2013 disclosures by Citigroup (\$1.5B FVA loss) and JPMorgan (\$1.5B FVA
charge), virtually all major dealers began reporting FVA. The ISDA produced a paper in 2014
documenting the range of industry approaches. The consensus methodology --- symmetric funding
curve applied to expected future funding gaps --- is what this model implements.

\subsection{Model Users}

\begin{itemize}[leftmargin=*]
  \item \textbf{Front Office Pricing}: FVA charge added to price of uncollateralised client derivatives.
  \item \textbf{Finance / P\&L}: Daily FVA P\&L recognised in income statement.
  \item \textbf{Treasury}: FVA informs funding cost allocation --- trading desks are charged for the funding they consume.
  \item \textbf{Collateral Negotiators}: FVA quantifies the value of upgrading a no-CSA relationship to a full CSA.
\end{itemize}

% ══════════════════════════════════════════════════════════════════════════════
\section{Theoretical Foundation}
% ══════════════════════════════════════════════════════════════════════════════

\subsection{FVA Formula}

The symmetric FVA formula under the risk-neutral measure is:

\begin{equation}
\mathrm{FVA} = -\int_0^T s_f(t)\,\mathbb{E}\!\left[V^+(t)\right] D(0,t)\,dt
             + \int_0^T s_f(t)\,\mathbb{E}\!\left[V^-(t)\right] D(0,t)\,dt
\tag{2}
\end{equation}

Decomposed into cost and benefit components:

\begin{align}
\mathrm{FCA} &= -\sum_i s_f \cdot \mathrm{ENE}_{\mathrm{coll}}(t_i) \cdot \Delta t_i \cdot D(0,t_i)
\tag{3} \\[6pt]
\mathrm{FBA} &= \phantom{-}\sum_i s_f \cdot \mathrm{EPE}_{\mathrm{coll}}(t_i) \cdot \Delta t_i \cdot D(0,t_i)
\tag{4}
\end{align}

where:
\begin{itemize}[leftmargin=*]
  \item $s_f$ = Apex's funding spread over OIS (current: 55\,bp for 1--5 year tenors)
  \item $\mathrm{EPE}_{\mathrm{coll}}(t)$ = expected positive exposure of the \emph{collateralised hedge} at time $t$
  \item $\mathrm{ENE}_{\mathrm{coll}}(t)$ = expected negative exposure of the collateralised hedge
  \item $D(0,t)$ = OIS discount factor
\end{itemize}

\textbf{Intuition}: FCA is paid when Apex receives money from the client trade but must fund
the hedge (borrowing cost). FBA is earned when Apex owes money on the client trade but can use
that receipt to fund itself more cheaply.

\subsection{Relationship to CVA}

CVA and FVA share the same simulation infrastructure (EPE/ENE profiles from Monte Carlo) but
differ in what drives the cost:

\begin{itemize}[leftmargin=*]
  \item \textbf{CVA}: the counterparty might \emph{default} (credit loss) --- proportional to default probability.
  \item \textbf{FVA}: Apex must \emph{fund} the position regardless of whether anyone defaults --- proportional to exposure profile $\times$ funding spread.
\end{itemize}

They are conceptually independent and both are priced into client derivatives quotes.

\subsection{The Asymmetric Variant (FCA Only)}

Some practitioners argue FBA should not be recognised because:
\begin{enumerate}[leftmargin=*]
  \item It requires Apex to be in a net payable position to a collateralised counterparty.
  \item The ``benefit'' of holding their cash cheaply is already priced in the OIS rate.
  \item Recognising FBA inflates P\&L on day 1 against a future cash flow that may not materialise.
\end{enumerate}

Apex's policy recognises symmetric FVA (FCA $-$ FBA) but requires the FBA component to be
separately disclosed in risk reports with a sensitivity showing the P\&L impact of removing FBA.

\subsection{Funding Curve Construction}

Apex's unsecured funding curve is constructed from:

\begin{enumerate}[leftmargin=*]
  \item \textbf{Short end (0--1Y)}: SOFR + Apex CP (commercial paper) spread from recent issuance.
  \item \textbf{Medium term (1--5Y)}: SOFR + Apex senior unsecured bond spread (Bloomberg/Markit secondary).
  \item \textbf{Long end (5--30Y)}: SOFR + Apex CDS spread + liquidity premium (estimated from basis between bond spread and CDS).
\end{enumerate}

The curve is rebuilt daily. Any point where Apex has not issued bonds in the past 12 months
uses linear interpolation between bracketing tenors.

\subsection{Literature References}

\begin{itemize}[leftmargin=*]
  \item Burgard, C., Kjaer, M.\ (2011). ``Partial Differential Equation Representations of Derivatives with Bilateral Counterparty Risk and Funding Costs.'' \textit{Journal of Credit Risk}, Vol.\ 7.
  \item Hull, J., White, A.\ (2012). ``The FVA Debate.'' \textit{Risk Magazine}, July 2012.
  \item Hull, J., White, A.\ (2014). ``Valuing Derivatives: Funding Value Adjustments and Fair Value.'' \textit{Financial Analysts Journal}, Vol.\ 70.
  \item Andersen, L., Duffie, D., Song, Y.\ (2019). ``Funding Value Adjustments.'' \textit{Journal of Finance}, Vol.\ 74.
  \item ISDA (2014). \textit{FVA: The ISDA Guidance}.
\end{itemize}

% ══════════════════════════════════════════════════════════════════════════════
\section{Data Requirements and Governance}
% ══════════════════════════════════════════════════════════════════════════════

\subsection{Input Data Sources}

\begin{table}[h!]
\centering
\caption{FVA Model Input Data Sources}
\begin{tabular}{@{}lll@{}}
\toprule
\textbf{Data Item} & \textbf{Source} & \textbf{Frequency} \\
\midrule
Apex funding curve (CP, bond spreads) & Treasury / Bloomberg & Daily \\
OIS discount curves (SOFR, ESTR, SONIA) & Bloomberg / Internal & Daily \\
EPE/ENE profiles & CVA model Monte Carlo engine & Daily (shared run) \\
Trade/netting set data & Front office systems & Real-time \\
CSA terms (VM threshold, MTA) & ISDA MA database & On amendment \\
\bottomrule
\end{tabular}
\end{table}

\subsection{Data Governance}

FVA shares the CVA model's simulation infrastructure. Any change to simulation parameters
(number of paths, time grid, market factor models) requires simultaneous review of both CVA and
FVA outputs to confirm consistency.

The funding curve construction is owned jointly by the XVA team and Treasury. Disputes about
the appropriate funding spread are escalated to the CFO / Head of Treasury.

% ══════════════════════════════════════════════════════════════════════════════
\section{Methodology and Implementation}
% ══════════════════════════════════════════════════════════════════════════════

\subsection{Shared Simulation with CVA}

FVA reuses the Monte Carlo simulation generated for CVA (same paths, same time steps, same
market factor evolution). The incremental computation is:
\begin{enumerate}[leftmargin=*]
  \item Apply netting set aggregation (same as CVA).
  \item Apply CSA collateral offset (same as CVA).
  \item Apply funding spread to EPE/ENE profiles.
\end{enumerate}

This integration reduces computational cost but creates a dependency: FVA results cannot be
produced without a successful CVA run.

\subsection{FVA for Partially Collateralised Trades}

Many CSAs have a threshold (e.g., Apex posts margin only when net MtM exceeds \$10M) or a
minimum transfer amount. These create a funding gap even for nominally collateralised
relationships:

\begin{equation}
\mathrm{Funding\_Gap}(t) = \max\!\bigl(V(t) - \mathrm{Threshold} - \mathrm{MTA},\,0\bigr) - C_{\mathrm{posted}}(t)
\tag{5}
\end{equation}

The FVA for partially collateralised trades is computed using this adjusted exposure rather
than the full EPE/ENE.

\subsection{Desk-Level FVA Allocation}

FVA is allocated to individual trading desks based on their contribution to the firm's aggregate
funding need. The allocation methodology:
\begin{enumerate}[leftmargin=*]
  \item Compute firm-level FVA.
  \item Decompose by netting set.
  \item Allocate netting set FVA to originating desk using trade-level sensitivities.
  \item Charge desk P\&L; credit is passed to Treasury (which provides the actual funding).
\end{enumerate}

This creates the correct incentive: traders who write large uncollateralised trades face
explicit P\&L charges for the funding they consume.

% ══════════════════════════════════════════════════════════════════════════════
\section{Model Testing and Backtesting}
% ══════════════════════════════════════════════════════════════════════════════

\subsection{Proxy Backtesting}

Daily FVA P\&L is compared against Greeks-predicted P\&L:

\begin{equation}
\Delta\mathrm{FVA}_{\mathrm{pred}} = \mathrm{DV01}_{\mathrm{FVA}} \cdot \Delta r
  + \Delta s_f \cdot \mathrm{Duration} \cdot \mathrm{Exposure}
\tag{6}
\end{equation}

Tolerance: P\&L explain $> 80\%$ of total FVA daily change. Lower threshold than CVA given the
additional uncertainty from funding spread moves.

\subsection{Benchmark Comparison}

Quarterly comparison against an independently computed FVA using a simplified closed-form for
a representative vanilla IRS:

\begin{equation}
\mathrm{FVA}_{\mathrm{approx}} \approx s_f \cdot \overline{\mathrm{EPE}} \cdot T
\tag{7}
\end{equation}

Agreement threshold: within 8\% of full simulation result.

\subsection{Sensitivity to Funding Curve}

Funding curve sensitivity analysis --- parallel shift of $+100\,\mathrm{bp}$ in Apex's funding
spread:
\begin{itemize}[leftmargin=*]
  \item Expected FVA increase: proportional to portfolio duration $\times$ exposure.
  \item Monitored quarterly; material sensitivity triggers review of hedging strategy.
\end{itemize}

% ══════════════════════════════════════════════════════════════════════════════
\section{Performance Metrics and Calibration Standards}
% ══════════════════════════════════════════════════════════════════════════════

\begin{table}[h!]
\centering
\caption{FVA Performance Metrics and Thresholds}
\begin{tabular}{@{}lll@{}}
\toprule
\textbf{Metric} & \textbf{Acceptable Range} & \textbf{Action if Breached} \\
\midrule
FVA P\&L Explain          & $>80\%$                & Investigate within 2 business days \\
Benchmark deviation       & $<8\%$ relative        & Review within 30 days \\
Funding curve staleness   & No tenor $>5$ days old & Data quality alert \\
FBA as \% of FCA          & Disclosed separately   & Quarterly disclosure to MRO \\
\bottomrule
\end{tabular}
\end{table}

\subsection{Recalibration Schedule}

\begin{itemize}[leftmargin=*]
  \item \textbf{Funding curve}: Daily (automated from bond/CDS feed).
  \item \textbf{Simulation paths (shared with CVA)}: No recalibration; path count reviewed annually.
  \item \textbf{FCA/FBA policy (asymmetric vs.\ symmetric)}: Annual review by CFO / CRO.
\end{itemize}

% ══════════════════════════════════════════════════════════════════════════════
\section{Limitations, Assumptions, and Known Weaknesses}
% ══════════════════════════════════════════════════════════════════════════════

\subsection{Theoretical Controversy}

The fundamental limitation of FVA is theoretical: the academic literature has not reached
consensus on whether FVA is a legitimate fair value adjustment or a distortion of asset pricing
principles. Hull and White's critique --- that FVA leads to negative NPV trades being accepted
when they should not be --- is mathematically valid in a frictionless market. In practice,
market frictions (unsecured funding is genuinely more expensive than collateralised) make FVA
operationally real regardless of its theoretical status.

\textbf{Compensating control}: FVA is disclosed separately from CVA and DVA in financial
statements. The sensitivity of P\&L to the FVA methodology (symmetric vs.\ FCA-only) is
disclosed in the quarterly risk report.

\subsection{Funding Spread Instability in Stress}

In a stress event affecting Apex's credit, the funding spread may widen dramatically and
non-linearly. The current model uses a flat funding spread calibrated to current market
conditions. A 100\,bp widening in Apex's credit spread would increase FCA by approximately
\$180M.

\textbf{Compensating control}: Stressed FVA computed using funding spread at the 75th
percentile of historical 6-month widening.

\subsection{Wrong-Way Funding Risk}

For some counterparties, the scenarios where Apex's funding costs are highest (stress events)
coincide with scenarios where those counterparties are most likely to increase their derivatives
exposure. This ``wrong-way funding risk'' is not explicitly modelled.

\textbf{Compensating control}: Concentration limits on uncollateralised exposure by desk ensure
no single counterparty drives excessive FVA.

% ══════════════════════════════════════════════════════════════════════════════
\section{Compensating Controls}
% ══════════════════════════════════════════════════════════════════════════════

\begin{table}[h!]
\centering
\caption{FVA Compensating Controls}
\begin{tabular}{@{}p{4.5cm}p{6cm}p{3.5cm}@{}}
\toprule
\textbf{Risk} & \textbf{Control} & \textbf{Owner} \\
\midrule
Theoretical controversy & Separate disclosure; FCA-only sensitivity & Finance / CFO \\[4pt]
Funding spread instability & Stressed FVA at 75th percentile spread & XVA Quant Team \\[4pt]
Funding curve staleness & Daily rebuild; 5-day staleness alert & Market Data \\[4pt]
FBA overstatement & FBA disclosed separately; policy review annually & CFO / CRO \\[4pt]
Shared simulation failure & Fallback: prior-day FVA + Greek adjustment & XVA Technology \\
\bottomrule
\end{tabular}
\end{table}

\subsection{Model Reserve}

A model reserve of \textbf{\$25 million} is held against FVA model uncertainty:

\begin{table}[h!]
\centering
\caption{FVA Model Reserve Breakdown}
\begin{tabular}{@{}lr@{}}
\toprule
\textbf{Reserve Component} & \textbf{Amount (USD M)} \\
\midrule
Theoretical methodology uncertainty (FCA vs.\ FCA--FBA) & 12 \\
Funding curve construction uncertainty                   & 8 \\
Wrong-way funding risk                                   & 5 \\
\midrule
\textbf{Total Reserve}                                   & \textbf{25} \\
\bottomrule
\end{tabular}
\end{table}

% ══════════════════════════════════════════════════════════════════════════════
\section{Change Management}
% ══════════════════════════════════════════════════════════════════════════════

Material changes (funding curve methodology, FCA/FBA policy, integration with new simulation
infrastructure) require full MRO re-validation. Non-material changes (new tenor point added to
funding curve, computational optimisation) proceed via expedited review with 5-business-day
notification to MRO.

\subsection{Material Change Triggers}

\begin{itemize}[leftmargin=*]
  \item Change to the definition of eligible hedging instruments (collateralised vs.\ uncollateralised).
  \item Change to the symmetric/asymmetric FCA/FBA policy.
  \item Replacement of the shared Monte Carlo simulation infrastructure.
  \item Change to the FTP (Funds Transfer Pricing) rate used as the funding curve basis.
\end{itemize}

% ══════════════════════════════════════════════════════════════════════════════
\section{Approvals and Sign-Off}
% ══════════════════════════════════════════════════════════════════════════════

\begin{table}[h!]
\centering
\caption{Approval Status}
\begin{tabular}{@{}llll@{}}
\toprule
\textbf{Role} & \textbf{Name} & \textbf{Status} & \textbf{Date} \\
\midrule
Model Developer          & XVA Quant Team    & Submitted & March 2026 \\
Model Risk Officer       & Dr.\ Rebecca Chen & Pending   & --- \\
Chief Risk Officer       & Dr.\ Priya Nair   & Pending   & --- \\
CFO (FVA Policy Owner)   & Diana Osei        & Pending   & --- \\
Business Owner           & James Okafor      & Pending   & --- \\
\bottomrule
\end{tabular}
\end{table}

\vspace{1cm}
\noindent\textcolor{apexgray}{\footnotesize
Document Classification: RESTRICTED --- Model Risk Sensitive\\
Apex Global Bank \textbar{} Quantitative Research Division \textbar{} XVA Team\\
Next scheduled review: March 2027 (or upon material model change)
}

\end{document}
